\section{Introduction}\label{sec:introduction}

A \emph{Cartan--Hadamard manifold} is a complete simply connected smooth
Riemannian manifold with nonpositive sectional curvature. The generalized
Cartan--Hadamard isoperimetric conjecture asserts that if its sectional
curvature is at most $\kappa\le0$, then every relatively compact smooth
domain has boundary area at least that of the equal-volume ball in the
simply connected space form of curvature $\kappa$. This compares arbitrary
domains, so geodesic-ball volume comparison alone does not suffice.
We prove the conjecture in every dimension, for arbitrary measurable sets
of finite volume and finite perimeter.

\subsection{The model comparison}

Write $\omega_n$ for the volume of the unit ball in $\R^n$, set
$m=n-1$, and put $s_m=|\mathbb S^m|=n\omega_n$. For $b\ge0$, define
\begin{equation}\label{eq:model-functions}
 \sn_b(r)=
 \begin{cases}r,&b=0,\\ b^{-1}\sinh(br),&b>0,\end{cases}
 \qquad \cs_b(r)=\sn_b'(r).
\end{equation}
The area and volume of a radius-$r$ ball in curvature $-b^2$ are
\begin{equation}\label{eq:model-area-volume}
 \mathcal A_b(r)=s_m\sn_b(r)^m,
 \qquad
 \mathcal V_b(r)=s_m\int_0^r\sn_b(t)^m\,dt.
\end{equation}
Both functions increase continuously from zero to infinity. The model
isoperimetric profile is
\[
 \mathcal I_b(V)=\mathcal A_b\bigl(\mathcal V_b^{-1}(V)\bigr),
 \qquad V\ge0.
\]
In particular, $\mathcal I_b(0)=0$, and $\mathcal I_b$ is continuous
and strictly increasing on $[0,\infty)$.

For a measurable set $E\subset M$, its \emph{ambient perimeter} is
\begin{equation}\label{eq:ambient-perimeter}
 P_g(E)=|D\chi_E|_g(M)
 =\sup\left\{\int_E\divg_g X\,d\vol_g:
 X\in C_c^\infty(TM),\ |X|_g\le1\right\}.
\end{equation}
Here $\chi_E$ is the indicator of $E$, and $|D\chi_E|_g$ is its
variation measure. A set has finite perimeter when this quantity is
finite. Sets are identified up to ambient volume zero. For a relatively
compact smooth domain, perimeter is its boundary area.

\begin{theorem}\label{thm:main}
Let $n\ge2$, let $\kappa\le0$, and let $(M^n,g)$ be a complete simply
connected smooth Riemannian manifold with $\Sec_g\le\kappa$.
Every measurable set $E\subset M$ with $\vol_g(E)<\infty$ and
$P_g(E)<\infty$ satisfies
\[
 P_g(E)\ge \mathcal I_{\sqrt{-\kappa}}\bigl(\vol_g(E)\bigr).
\]
\end{theorem}

For $\kappa=0$, the conclusion reads
\[
 P_g(E)
 \ge n\omega_n^{1/n}\vol_g(E)^{(n-1)/n}.
\]
For $\kappa=-1$, if
$\vol_g(E)=n\omega_n\int_0^r\sinh^{n-1}t\,dt$, it reads
\[
 P_g(E)\ge n\omega_n\sinh^{n-1}r.
\]
The comparison concerns full ambient perimeter, including boundary on
any confining sphere used in the proof. It applies to unbounded sets
of finite volume as well as bounded ones. Balls in the model spaces
attain the bounds.

The Euclidean equality case has the following classification.

\begin{theorem}[Euclidean equality rigidity]\label{thm:equality}
Let $(M^n,g)$ be complete, simply connected and smooth, with
$n\ge2$ and $\Sec_g\le0$. Let $E\subset M$ be bounded and measurable,
with $0<\vol_g(E)<\infty$ and $P_g(E)<\infty$. Then
\[
 P_g(E)=n\omega_n^{1/n}\vol_g(E)^{(n-1)/n}
\]
if and only if $E$ agrees up to ambient volume zero with an open region
whose induced Riemannian metric is isometric to a round Euclidean ball.
\end{theorem}

Thus rigidity determines the metric on the entire equality region.
It does not require the ambient manifold outside that region to be
flat. Boundedness and positive volume belong to the equality theorem;
Theorem~\ref{thm:main} has the full finite-volume scope stated above.

\subsection{Model-ball spectral and torsional consequences}

For a smooth relatively compact domain $D$ in a Riemannian manifold
$(M,g)$ and $1<p<\infty$, let
$W^{1,p}_0(D)$ be the closure of $C_c^\infty(D)$ in the Sobolev norm, and
define
\[
 \begin{aligned}
 \lambda_{1,p}(D)&=\inf_{0\ne u\in W^{1,p}_0(D)}
 \frac{\int_D|\nabla_g u|_g^p\,d\vol_g}{\int_D|u|^p\,d\vol_g},\\
 T_p(D)&=\sup_{0\ne u\in W^{1,p}_0(D)}
 \frac{\bigl(\int_D|u|\,d\vol_g\bigr)^p}
      {\int_D|\nabla_g u|_g^p\,d\vol_g}.
 \end{aligned}
\]
The first is the Dirichlet variational value for
$-\Delta_{p,g}u=\lambda|u|^{p-2}u$, where\newline
$\Delta_{p,g}u=\divg_g(|\nabla_g u|_g^{p-2}\nabla_g u)$ is the
unnormalized $p$-Laplacian. The second uses the quotient normalization
of $p$-torsional rigidity in \cite[(1.1)--(1.6)]{BrianiButtazzoPrinari2022}.

\begin{corollary}[Model-ball Faber--Krahn and Saint--Venant comparisons]
\label{cor:model-ball-spectral}
Let $n\ge2$, $\kappa\le0$ and $1<p<\infty$. Let $(M^n,g)$ be complete,
simply connected, smooth and without boundary, with $\Sec_g\le\kappa$.
For a nonempty smooth domain $D\Subset M$, let $B_\kappa(|D|)$ be a
geodesic ball of volume $|D|=\vol_g(D)$ in the simply connected
$n$-dimensional space form of curvature $\kappa$. Then
\[
 \lambda_{1,p}(D)\ge\lambda_{1,p}\bigl(B_\kappa(|D|)\bigr),
 \qquad T_p(D)\le T_p\bigl(B_\kappa(|D|)\bigr).
\]
\end{corollary}

Theorem~\ref{thm:main} supplies the isoperimetric input for
equimeasurable radial rearrangement into the model ball
\cite{NobiliViolo2025}. This rearrangement preserves the $L^1$ and
$L^p$ integrals and does not increase the $p$-Dirichlet energy,
giving the two quotient comparisons. The proof appears in
Section~\ref{sec:consequences}.

\subsection{History and related work}

The surface inequality originates in Weil's work and the
subharmonic-function approach of Beckenbach and Rad\'o; Bol developed
the comparison with nonzero model curvature
\cite{Weil1926,BeckenbachRado1933,Bol1941}.
We use the smooth surface formulation of Chavel and Feldman
\cite{ChavelFeldman1980}. Kleiner proved the model comparison in
dimension three, including its equality case
\cite[Theorem~2]{Kleiner1992}. Croke proved the sharp Euclidean
comparison in dimension four \cite{Croke1984}. Kloeckner and Kuperberg
subsequently developed optical and chord-based comparison methods in
dimensions two and four; their four-dimensional negative-curvature
comparison includes an explicit restriction on the domain's chord
length and the radius of its model ball
\cite[Theorem~1.5]{KloecknerKuperberg2019}.

A second line of work connects isoperimetry to total curvature.
Ghomi and Spruck showed that, in a fixed Cartan--Hadamard
$n$-manifold, the lower bound $\int_\Gamma|\mathrm{GK}|\ge s_{n-1}$
for every closed convex $C^{1,1}$ hypersurface $\Gamma$ implies the
Euclidean isoperimetric inequality for bounded positive-volume
finite-perimeter sets, with equality only for regions isometric to
Euclidean balls \cite[Theorem~7.1]{GhomiSpruck2022}. Here
$\mathrm{GK}$ is the Gauss--Kronecker curvature, and such a convex
hypersurface bounds a compact convex set with nonempty interior. Their analysis of minimizers
confined to a geodesic ball includes the obstacle regularity that we
use \cite[Lemma~7.2(i),(ii)]{GhomiSpruck2022}; this regularity lemma
does not assume the total-curvature inequality. Ghomi and Stavroulakis
prove an inequality for sufficiently small $C^2$ perturbations of a
Euclidean ball metric with nonpositive curvature
\cite{GhomiStavroulakis2026}. Their theorem concerns the metric on the
ball itself and does not require an extension to a complete
Cartan--Hadamard manifold. Related results also include Ghomi's
inequalities under negative-curvature pinching \cite{Ghomi2026Pinched}.

Recent preprints overlap with several parts of the problem.
Chen, Ghomi and Wang establish Euclidean comparison and rigidity in
dimensions three through nine, report calibration verification through
at least dimension thirteen, and announce generalized comparison in
dimensions four and five
\cite[Theorems~1.1--1.2 and the subsequent discussion]{ChenGhomiWang2026HigherDimensions}.
Their argument uses boundary-pair integrals, Jacobi fields and a
confined profile; its appendix includes exact-arithmetic certificates.
Agnoletto, Da Silva, Nardulli and Resende prove the small-volume
model comparison under a Ricci lower bound, and reduce the
unrestricted-volume comparison to equality rigidity in a
noncollapsing boundaryless Alexandrov class with two-sided curvature
bounds \cite[Theorems~1.1 and 1.3, Assumption~1]{AgnolettoDaSilvaNardulliResende2026}.
Wheeler proves all-dimensional Euclidean comparison under controlled
variation of a variable negative-curvature scale
\cite[Theorem~1.1 and Definition~1.2]{Wheeler2026}.
These dimensional or curvature restrictions, and the Alexandrov
rigidity condition in the reduction, differ from the hypotheses of
Theorems~\ref{thm:main} and~\ref{thm:equality}.

\subsection{Main ideas and proof organization}

The central analytic estimate is a sharp critical Sobolev inequality
on a hypersurface. Normalize the curvature bound to $-b^2$, with
$b\in\{0,1\}$. For a locally $C^{1,1}$ immersed hypersurface
$\Sigma^m$ with $m=n-1\ge3$, let $d\sigma$ be its induced area
measure, $\nabla_\Sigma$ its tangential gradient, and $h$ its signed
normalized scalar mean curvature for a local unit normal. Put
\[
 q=\frac{2m}{m-2},\qquad c=\frac4{m-2},\qquad Y_m=m s_m^{2/m}.
\]
We prove
\begin{equation}\label{eq:intro-sobolev}
 \int_\Sigma\left(c|\nabla_\Sigma v|^2+m(h^2-b^2)v^2\right)\,d\sigma
 \ge Y_m\left(\int_\Sigma|v|^q\,d\sigma\right)^{2/q}
\end{equation}
for Lipschitz tests compactly supported in the interior of $\Sigma$.
The square $h^2$ is independent of the local normal. The hypersurface
need not be complete or connected. The coefficient is the sharp
Euclidean critical Sobolev constant of Aubin and Talenti in this
normalization \cite{Aubin1976,Talenti1976}.

Sections~\ref{sec:maps} and~\ref{sec:sobolev} establish this estimate.
Radial maps into a unit sphere come with positive weights whose
logarithmic Hessians and map differentials satisfy compatible
curvature-comparison bounds. The center and scale can be chosen so
that every spherical coordinate has zero mean for a prescribed
compactly supported nonatomic probability measure. The bounds combine
into a completed square in the hypersurface energy.
If the critical quotient were below the Euclidean threshold, local
compactness and Br\'ezis--Lieb splitting would give a Dirichlet
minimizer \cite{BrezisLieb1983}. Balancing its critical mass and using
the spherical coordinates as variations forces a first-order identity.
A weighted zero extension then has zero gradient and nonzero mass,
contradicting the Dirichlet condition. The centering and coordinate
variations are related to Li and Yau's conformal-volume argument
\cite[proof of Theorem~1]{LiYau1982}; the maps and weighted centering
needed here are constructed directly.

To compare domains, the difficulty is that the mean curvature of an
arbitrary boundary is uncontrolled. Sections~\ref{sec:profile}
and~\ref{sec:comparison} instead minimize full ambient perimeter at
fixed volume inside a large geodesic ball. This is the confined-profile
strategy of Kleiner, with the regularity formulation of Ghomi and Spruck
\cite{Kleiner1992,GhomiSpruck2022}. Free variations identify the
boundary curvature with the profile derivative; inward variations
control it at contact with the confining sphere. A deletion-and-volume-repair argument first proves area growth near the singular set.
Cutoffs with vanishing Dirichlet energy then make the constant test
legitimate in \eqref{eq:intro-sobolev}. This gives one differential
inequality whose integral is exactly the model comparison.
Section~\ref{sec:dimensions} supplies the classical two- and
three-dimensional inputs, their strict BV approximation, a common
finite-volume cutoff argument, and metric scaling to arbitrary
$\kappa<0$.

For equality, Section~\ref{sec:equality-regularity} first uses the
proved Euclidean comparison to make a bounded equality set a local
perimeter almost-minimizer. Local regularity
\cite[Corollary~1.6]{AntonelliPasqualettoPozzetta2022}, followed by a
difference-quotient argument, gives a locally $C^{1,1}$ regular
boundary with the same signed constant mean curvature on every
component. Singular cutoffs extend first variation and the Sobolev
inequality to ambient tests on the whole perimeter support.

Section~\ref{sec:rigidity} locates that support and then recovers the
region. In dimensions at least four, the constant equality test in
\eqref{eq:intro-sobolev} gives a sharp spectral bound. Logarithm
coordinates centered at a squared-distance barycenter use this bound
to give an upper estimate for the radial second moment. A radial flux
identity gives the matching lower estimate. Equality forces the
outward normal to be radial with constant radius, so the perimeter
support lies on a sphere. Periodic Wirtinger replaces the spectral
argument in dimension two; a punctured flux identity replaces it in
dimension three and produces a sphere in a tangent space whose center
may be offset from the logarithm origin. In all dimensions, BV phase
constancy and boundedness identify the occupied interior. Exponential
metric domination and equality of volume then force every metric
eigenvalue to be Euclidean throughout that ball. This last step proves
whole-region rigidity, beyond the preceding containment of its boundary.

\subsection{Conventions}

All hypersurface metrics and area measures are induced. For a local
unit normal $N$, we use the signed trace curvature
$H=\divg_\Sigma N$ and the signed normalized scalar $h=H/m$.
For a smooth ambient function $u$, its restriction satisfies
\[
 \Delta_\Sigma u=\tr_{T\Sigma}\Hess_M u-HNu.
\]
Thus a Euclidean sphere of radius $r$ has outward $h=1/r$, and a
curvature-$-1$ sphere has $h=\coth r$. Reversing $N$ reverses $h$;
only normal-independent expressions are used when no global normal is
chosen. For boundaries of sets, the occupied phase fixes $N$.
All functions are real-valued. We often write $|E|=\vol_g(E)$ and
$P(E)=P_g(E)$, and use $P(E;D)=|D\chi_E|(D)$ for the perimeter
measure on a Borel set $D$. For $p\in M$ we write
$\log_p=\exp_p^{-1}$, $r_p(x)=d(p,x)$ and $D_p=r_p^2/2$;
Section~\ref{sec:maps} proves the comparisons used for these functions.
